% Unidad R566: introducción, líneas francesas 24--35, pp. 5--6.
% Traducción: OpenAI Codex, GPT-6.1 Sol, esfuerzo Ultra.
No corresponde en esta Introducción ofrecer una descripción más o menos sumaria del punto de vista de los «esquemas» en geometría algebraica, ni la larga lista de razones que hicieron necesaria su adopción, en particular la aceptación sistemática de elementos nilpotentes en los anillos locales de las «variedades» que consideramos (lo que necesariamente relega a un segundo plano la noción de aplicación racional, en favor de la de aplicación regular o «morfismo»).
Este tratado se propone precisamente desarrollar de manera sistemática el lenguaje de los «esquemas» y demostrará, así lo esperamos, su necesidad.
Aunque sería fácil hacerlo, tampoco intentaremos ofrecer aquí

\oldpage[I]{6}
una introducción «intuitiva» a las nociones desarrolladas en el capítulo primero.
El lector que desee tener una visión preliminar de las materias de este tratado podrá consultar la conferencia de A.~Grothendieck en el Congreso Internacional de Matemáticos celebrado en Edimburgo en 1958~\cite{I-7}, así como la exposición~\cite{I-8} del mismo autor.
El trabajo~\cite{I-14} (citado como~(FAC)) de J.-P.~Serre también puede considerarse una exposición intermedia entre el punto de vista clásico y el de los esquemas en geometría algebraica; por ello, su lectura puede constituir una excelente preparación para la de nuestros \emph{Éléments}.

\bigskip
\begin{center}*\quad*\quad*\end{center}
\bigskip
